对随机样本的定义有两点特别值得注意:

1.\quad 在单个试验中, $p$ 个变量的测量结果, 例如 $\mathbf{X}_{j}'=[X_{j1},X_{j2},\cdots,X_{jp}]$, 常常是相关的. 的确, 我们期望这是事实. 然而, 从不同试验得到的测量结果必须是独立的.

2.\quad 当变量可能会随时间缓慢地变化时, 例如 $p$ 种股票价格或 $p$ 个经济指标构成的集合, 各次试验的测量结果的独立性可能不成立. 违背独立性的临时假设对统计推断的性质会有严重影响.

下列例子将说明这些道理.

\medskip
\noindent\textbf{例 3.5}\quad (选择一个随机样本)

为了设计一种准入制度使野外划舟区不至于太拥挤, 一位自然资源经理首先对客户进行了调查. 整个野外区域被分成许多子区域, 并要求应答者给出有关访问区域、逗留时间及其他变量的信息.

接下来的办法是, 在特定一周内对进入该野外区域的所有人进行随机抽选 (可以利用随机数表). 所有人员以等可能的机会被选为样本, 因此, 划舟者的比例越大表示进入此区域的人越多.

这里, 我们会期望样本观测值严格符合用户或潜在用户总体的随机样本的标准. 另一方面, 如果一个取样者等待在深入区域内部的营地, 并仅调查到达该地的划舟者, 那么相继的观测结果不会是独立的. 例如, 停留时间会趋向于较长.\hfill$\blacksquare$

\medskip
\noindent\textbf{例 3.6}\quad (非随机样本)

由于关注未来如何对固体废品进行处理, 因而对美国每年产生的固体废品的总重量进行了研究 (环境保护署). \hyperref[tab:3.1]{表 3.1} 给出了所选年份中 $x_{1}=$ 废纸及废记录纸和 $x_{2}=$ 塑料废品 (按百万吨计算) 的估计总量. $\mathbf{X}'=[X_{1},X_{2}]$ 中的这些观测结果可否作为一个 $n=7$ 的随机样本来处理? 不! 其实除了在 2003 年的废纸及废纸板数据略有下降外, 这两个变量是随时间而增加的. 各年的值是来自同一分布的独立观测值, 这种情况是很少的.

\begin{mvtable}
\phantomsection\label{tab:3.1}\mvtabcap{表 3.1\quad 固体废品}
\centering
\begin{tabular}{lccccccc}
\toprule
年 & 1960 & 1970 & 1980 & 1990 & 1995 & 2000 & 2003\\
\midrule
$x_{1}$(纸) & 29.2 & 44.3 & 55.2 & 72.7 & 81.7 & 87.7 & 83.1\\
$x_{2}$(塑料) & 0.4 & 2.9 & 6.8 & 17.1 & 18.9 & 24.7 & 26.7\\
\bottomrule
\end{tabular}
\end{mvtable}

在第 1 章中, 我们初步讨论过, 统计独立的概念对测量距离有重要的意义. 若向量的分量之间是独立的并有相同的方差, 欧几里得距离是合适的. 假定我们考虑作为 $n$ 维空间中一点 $\mathbf{X}$ 的第 $k$ 列 $\mathbf{Y}_{k}=[X_{1k},X_{2k},\cdots,X_{nk}]$ 的位置. 该点的位置是由联合分布 $f(\mathbf{y}_{k})=f(x_{1k},x_{2k},\cdots,x_{nk})$ 确定的. 当测量结果 $X_{1k},X_{2k},\cdots,X_{nk}$ 是一随机样本时, $f(\mathbf{y}_{k})=f(x_{1k},x_{2k},\cdots,x_{nk})=f_{k}(x_{1k})f_{k}(x_{2k})\cdots f_{k}(x_{nk})$. 因此, 每个坐标 $x_{ik}$ 通过相同的边缘分布 $f_{k}(x_{jk})$ 同等地影响着位置.

若 $n$ 个分量不是独立的或边缘分布是不相同的, 单个的观测结果 (坐标) 对位置的影响是不对称的. 因此, 应考虑一个距离函数, 其中不同的坐标将不等地加权, 如同第 1 和第 2 章中介绍的“统计”距离或二次型.

在不对变量的基本联合分布的形式作更多假定的情况下, 可以得到关于 $\bar{\mathbf{x}}$ 和 $\mathbf{S}_{n}$ 的抽样分布的某些结论. 具体地说, 我们可以看到 $\bar{\mathbf{x}}$ 和 $\mathbf{S}_{n}$ 如何作为相应的总体均值向量 $\boldsymbol{\mu}$ 和协方差矩阵 $\boldsymbol{\Sigma}$ 的点估计.
